Analyze the ARCH(1) model
An ARCH(1) process \(\varepsilon\) satisfies the following equations:

\(\varepsilon_t = \sigma_t Z_t\)
\(\sigma_t^2 = \alpha_0 + \alpha_1 \varepsilon_{t - 1}^2\)

Weak stationarity requires \(\alpha_1 < 1\).

Unconditional variance: The unconditional variance is
\(\text{Var}(\varepsilon_t) = \sigma^2 = \frac{\alpha_0}{1 - \alpha_1}\)
Leptokurtic distribution: Suppose that \(Z_t \sim N(0, 1)\). Then for \(3 \alpha_1^2 < 1\):
\(\text{Kurt}(\varepsilon_t) = \frac{E(\varepsilon_t^4)}{E(\varepsilon_t^2)} = 3 \cdot \frac{1 - \alpha_1^2}{1 - 3 \alpha_1^2} \ge 3,\)
which is larger than \( \text{Kurt}(N) \) with \( N \sim \mathcal{N}(0, 1) \), so we have fat tails.


Consider an ARCH(1) process \(\varepsilon\) with finite fourth moments and \(Z_t \sim N(0, 1)\).

1. The process admits the following representation convergent in \(L^2\):
\(\varepsilon_t = \alpha_0 \sum_{k = 0}^\infty \alpha_1^k \prod_{j = 0}^k Z_{t - j}^2\)
2. \(\varepsilon^2\) is an AR(1) process with
\(\varepsilon_t^2 = \alpha_0 + \alpha_1 \varepsilon_{t - 1}^2 + \eta_t\)
and white noise \(\eta_t = \sigma_t^2 (Z_t^2 - 1)\).

Some sufficient conditions strict stationarity of an ARCH(1) are
\( \mathbb{E}[\log(α_1 Z_t^2 )]< 0\), which gives
\( α_1 < 3.5622 \) for \( Z_t \sim \mathcal{N}(0, 1) \).
Covariance stationarity, i.e., the \( L^2 \)-condition, requires \( \alpha_1 < 1\)


Define GARCH processes and state their representation


Let \((Z_t)_{t \in \mathbb{Z}}\) be SWN(0, 1). 
A process \((\varepsilon_t)_{t \in \mathbb{Z}}\) is called GARCH(\(p, q\)) process 
(generalized autoregressive conditionally heteroscedastic), if:

1. \((\varepsilon_t)_{t \in \mathbb{Z}}\) is strictly stationary.
2. For some strictly positive process \((\sigma_t)_{t \in \mathbb{Z}}\) the following identities hold:
with
\(\varepsilon_t = \sigma_t Z_t\)
\(\sigma_t^2 = \alpha_0 + \sum_{i = 1}^p \alpha_i \varepsilon_{t - i}^2 + \sum_{j = 1}^q \beta_j \sigma_{t - j}^2\)
with constants \(\alpha_0 > 0, \alpha_1, ..., \alpha_p, \beta_1, ..., \beta_q \ge 0\).

Similar to the ARCH process we can derive

Conditional expectation \( \mathbb{E}[\varepsilon_t\mid \mathcal{F}_{t-1}] = 0 \)
so \( \varepsilon \) is a martingale difference sequence,
however for the volatility
\( \sqrt{\text{Var}[\varepsilon_t \mid \mathcal{F}_{t-1}} = \sigma_t \)
which implies
Heteroscedacity: Volatility / conditional variance vary randomly over time.
Volatility clusters: Periods of high volatility are more persistent than for ARCH as \( \sigma^2_t \) 
depends on \( \sigma^2_s \) for \( s \leq t \).


Let \(\epsilon\) be a strong GARCH(\(p, q\)) process with finite fourth moments. 
Then the stochastic process \(\epsilon^2\) can be represented as an ARMA(\(m, q\)) process with \(m = \max(p, q)\):

\(\epsilon_t^2 = \alpha_0 + \sum_{i = 1}^m \gamma_i \epsilon_{t - i}^2 - \sum_{j = 1}^p \beta_j \eta_{t - j} + \eta_t\)

with white noise \(\eta_t = \sigma_t^2 (Z_t^2 - 1)\) and \(\gamma_i = \alpha_i + \beta_i\).

The white noise \(\eta\) is conditionally heteroscedastic, i.e., it is not strict white noise.
The methods available for ARMA processes can be exploited for the analysis of GARCH processes.
Beyond our previous discussion on estimation, we will also discuss likelihood methods.


Describe how likelihood methods can be applied for the estimation of a GARCH(1,1) process.


The multivariate density can be expressed in terms of conditional densities:
\(
f_{\varepsilon_0, \varepsilon_1, \dots, \varepsilon_n}(x_0, x_1, \dots, x_n) = 
f_{\varepsilon_0}(x_0) \prod_{t=1}^{n} f_{\varepsilon_t | \varepsilon_{t-1}, \dots, \varepsilon_0} (x_t | x_{t-1}, \dots, x_0).
\)
The identity \( \varepsilon_t = \sigma_t Z_t \) and the independence of \( Z_t \) and \( \varepsilon_{t-1}, \dots, \varepsilon_0 \) implies
\(
f_{\varepsilon_t | \varepsilon_{t-1}, \dots, \varepsilon_0} (x_t | x_{t-1}, \dots, x_0) = g\left( \frac{x_t}{\sigma_t} \right) \frac{1}{\sigma_t}
\)
with density \( g \) of \( Z_t \) and \( \sigma_t = \sqrt{\alpha_0 + \alpha_1 x_{t-1}^2 + \beta_1 \sigma_{t-1}^2} \) 
(note how this recursive structure allows us to determine all \( \sigma_t \), if \( \sigma_0, \alpha_i \) is known)
In summary, we obtain
\(
f_{\varepsilon_0, \varepsilon_1, \dots, \varepsilon_n}(x_0, x_1, \dots, x_n) = f_{\varepsilon_0}(x_0) \prod_{t=1}^{n} g\left( \frac{x_t}{\sigma_t} \right) \frac{1}{\sigma_t}.
\)

Unfortunately, the density \( f_{\varepsilon_0} \) and \( \sigma_0^2 \) are unknown. As a simplification, 
one approximates the likelihood by the conditional density given \( \varepsilon_0 \) and chooses a fixed value for \( \sigma_0^2 \), 
e.g. the sample variance of \( x_0, x_1, \dots, x_n \). 
This provides an approximate likelihood function which is used for estimation:
\(
L(\alpha_0, \alpha_1, \beta_1, x_0, x_1, \dots, x_n) := f_{\varepsilon_1, \dots, \varepsilon_n | \varepsilon_0} (x_1, \dots, x_n) = \prod_{t=1}^{n} g\left( \frac{x_t}{\sigma_t} \right) \frac{1}{\sigma_t}.
\)

Due to the multiplicative structure of the approximated likelihood, the log-likelihood function is more convenient:
\(
\log L(\alpha_0, \alpha_1, \beta_1, x_0, x_1, \dots, x_n) = \sum_{t=1}^{n} \log \left\{ g\left( \frac{x_t}{\sigma_t} \right) \frac{1}{\sigma_t} \right\}.
\)

The maximum likelihood estimators can be numerically determined from the first order conditions:
\(
\frac{\partial}{\partial \alpha_0} \log L = 0,\frac{\partial}{\partial \alpha_1} \log L = 0,\frac{\partial}{\partial \beta_1} \log L = 0
\)


How can conditional risk measurements based on coherent monetary risk measures
be computed in GARCH models?


Let \( V_t \) be the value of a portfolio at time \( t \geq 0 \).  
The time series of the portfolio value induces a sequence of returns with flipped sign:
\(
\varepsilon_t := -\frac{V_t - V_{t-1}}{V_{t-1}}.
\)

Assume that this time series can be described by a GARCH(p,q) model \( \varepsilon_t = \sigma_t Z_t \). The estimation is based on historical data.  
On the basis of this approach, the losses in period \([t, t+1]\) are
\(
L_{t+1} = -(V_{t+1} - V_t) = V_t \varepsilon_{t+1} = V_t \sigma_{t+1} Z_{t+1}.
\)

Conditional loss distribution:
\(
\mathbb{P}[L_{t+1} \leq l \mid \mathcal{F}_t] = \mathbb{P}[Z_{t+1} \leq \frac{l}{V_t \sigma_{t+1}} \mid \mathcal{F}_t] = G\left( \frac{l}{V_t \sigma_{t+1}} \right)
\)
The distribution function of \( Z_{t+1} \) is denoted by \( G \).
Conditional Value at Risk:
\(
\operatorname{V@R}^{t}_{\alpha} (L_{t+1}) := \inf \{ l \in \mathbb{R} \mid \mathbb{P}[L_{t+1} > l \mid \mathcal{F}_t] \leq 1 - \alpha \}
\)
\(
= V_t \sigma_{t+1} \operatorname{V@R}^{t}_{\alpha} (Z_{t+1}) = V_t \sigma_{t+1} G^{(-1)} (\alpha)
\)
Conditional Average Value at Risk:
\(
\operatorname{AV@R}^{t}_{\alpha} (L_{t+1}) := \frac{1}{1 - \alpha} \int_{\alpha}^{1} \operatorname{V@R}^{t}_{u} (L_{t+1}) \, du = V_t \sigma_{t+1} \frac{1}{1 - \alpha} \int_{\alpha}^{1} G^{(-1)} (u) \, du
\)

Explain the motivation for EGARCH and TARCH and provide their definitions.

Volatility increases with the size of shocks.
But there is also an empirically documented asymmetry between positive and negative shocks. Negative
shocks lead to larger volatility than positive shocks.
A potential rational for this observation, the leverage effect.
A resolution is available via the EGARCH, TARCH models:

EGARCH
\(
\log \sigma_t^2 = w_t + \sum_{k=1}^{\infty} \beta_k g(Z_{t-k}), \quad g(Z_t) = \theta Z_t + \gamma |Z_t|
\)

with deterministic coefficients \( w_t, \beta_k, \theta, \gamma \) and strict white noise \( Z = (Z_t) \) with expectation 0 and variance 1.
Sign effect:  
The parameter \( \theta \) is typically negative, and the term \( \theta Z_t \) increases volatility the more, the smaller \( Z_t \).
Size effect:  
The parameter \( \gamma \) is typically positive, and the term \( \gamma |Z_t| \) increases volatility proportionally 
to the absolute value of the shocks.

Threshold ARCH (TARCH) simply checks, if the shocks take negative values using an indicator function. 
Apart from this, the dependence on the past returns or squared returns is simply linear:
\(
\sigma_t^{\delta} = w + \sum_{i=1}^{q} \alpha_i \varepsilon_{t-i}^{\delta} + \sum_{i=1}^{q} \alpha_i^{-} \varepsilon_{t-i}^{\delta} \mathbb{1}_{\{\varepsilon_{t-i}<0\}}
\)
with deterministic parameters \( w, \alpha_i, \alpha_i^{-} \) and \( \delta = 1,2 \).
The model can be extended by adding lagged conditional variances or standard deviations as regressors; such a model is called TGARCH.




Define copulas


A d-dimensional copula \(C : [0, 1]^d \rightarrow [0, 1]\) is the distribution function of a d-dimensional random vector 
with uniform one-dimensional marginal distributions \(U([0,1])\).

A copula \(C\) can equivalently be characterized by the following properties:
a) \(C(u_1, \ldots, u_d)\) is increasing in each component \(u_i\).
b) For all \(i \in \{1, \ldots, d\}\), \(u_i \in [0, 1]\): \(C(1, \ldots, 1, u_i, 1, \ldots, 1) = u_i\).
c) For all \((a_1, \ldots, a_d), (b_1, \ldots, b_d) \in [0, 1]^d\) with \(a_i \leq b_i, u_{j,1} = a_j, u_{j,2} = b_j, j = 1, 2, \ldots, d\), the following inequality holds:
\(
\sum_{i_1=1}^2 \cdots \sum_{i_d=1}^2 (-1)^{i_1 + \cdots + i_d} C(u_{1, i_1}, \ldots, u_{d, i_d}) \geq 0.
\)


(a) This property must hold, since a copula is a multivariate distribution function.
(b) Marginal distributions are uniform — implying this property.
(c) This condition is less obvious, but can be motivated as follows:
Let \((U_1, \ldots, U_n)^\top\) be a random vector with distribution function \(C\). Then
\(
\mathbb{P}[U_1 \in (a_1, b_1], \ldots, U_d \in (a_d, b_d]] \geq 0.
\)
Computing this probability in terms of the distribution function \(C\) implies condition c).


State Sklar's Theorem

1. 
For any \(d\)-dimensional distribution function \(F\) with margins \(F_1, \ldots, F_d\) there exists a copula \(C\) with
\(
F(x_1, \ldots, x_d) = C(F_1(x_1), \ldots, F_d(x_d)) \quad \text{für alle } x_1, \ldots, x_d \in [-\infty, \infty].
\)
2.
Are the marginal distribution functions \(F_1, \ldots, F_d\) continuous, then the copula is uniquely given by
\(
C(u_1, \ldots, u_d) = F(F_1^{-1}(u_1), \ldots, F_d^{-1}(u_d)).
\)
Otherwise the copula is uniquely specified on \(\operatorname{Ran} F_1 \times \ldots \times \operatorname{Ran} F_d\), where 
\(\operatorname{Ran} F_i\) denotes the range of \(F_i\), \(i = 1, 2, \ldots, d\)
3.
Conversely, for a given copula \(C\) and given one-dimensional distribution functions \(F_1, \ldots, F_d\), 
the function \(F\) in (1) a \(d\)-dimensional distribution function with copula \(C\) and marginal distribution functions \(F_1, \ldots, F_d\).

The interpretationas of each statement are:
(1)
The marginal distributions describe the univariate risks.
Stochastic dependencies between risks are encoded by the copula.
(2)
\(F_i^{(-1)}(u_i)\) signifies the \(u_i\)-quantile of \(X_i\) or, equivalently, the distribution function \(F_i\).
\(C(u_1, \ldots, u_d)\) is the probability that the random variables \(X_1, \ldots, X_d\) do not exceed their quantiles at levels \(u_1, \ldots, u_d\).
A copula expresses dependencies on the level of quantiles.
Application: Using 2. we can construct implicit copulas on the basis of multivariate distributions with continuous margins.
(3)
Multivariate distributions may be specified on the basis of marginal distributions and a copula.
In practice, marginal distributions are often better understood than dependence structures. 
Copulas and Sklar's characterization of multivariate distributions admits a combination of marginal distributions and different copulas; 
the influence of different types of dependence can be analyzed in these models.

Describe the basic construction to create copulas (specializations come later)


Specify the marginal distributions \(F_1, \ldots, F_d\).
Choose a copula \(C\) and set \(F(x_1, \ldots, x_d) = C(F_1(x_1), \ldots, F_d(x_d)).\)

For each choice of \(C\) the multivariate distribution function \(F\) defines a distribution with the same univariate margins.
A given set of different copulas defines a class of such multivariate distributions.

For a detailed analysis of the problem, select appropriate copulas \(C\) that reflect the available information on dependence.

Another step is the monotonic transformations of the margins
We focus on continuous marginal distributions \(F_1, \ldots, F_d\)

The unique copula associated with a random vector \(\mathbf{X} = (X_1, \ldots, X_d)^\top\) with distribution function \(F\) is 
\(C(u_1, \ldots, u_d) = F(F_1^{(-1)}(u_1), \ldots, F_d^{(-1)}(u_d)) = \mathbb{P}[F(X_1) \leq u_1, \ldots, F(X_d) \leq u_d].\)

The copula \( C \) is invariant under strictly monotonic transformations of the margins, i.e., 
suppose that \(T_1, \ldots T_d\) are strictly increasing, then the random vector \( (T_1(X_1), \ldots, T_d(X_d))^\top \) 
has the same copula \( C \) as the random vector \( (X_1, \ldots, X_d)^\top \).


State the Frechet bounds for Copula's

Any copula \(C\) is bounded from above and below by the Fréchet bounds:
\(\max \left( \sum_{i = 1}^d u_i + 1 - d, 0 \right) \le C(u) \le \min(u_1, u_2, ..., u_d)\)

with \(u = (u_1, u_2, ..., u_d)\).
* The lower bound is a copula only in the case \(d = 2\), and is called countermonotonic copula in this case.
* The upper bound is the comonotonic copula.

In fact this holds for any multivariate distribution \( F \) with margins \( F_i \) (replace \( C(u) \) by \( F \) and \( u_i \) by \( F_i \)).


Give the heuristic categories of Copula's and describe the fundamental ones.

Fundamental Copulas:
Copulas associated with simple and important types of dependence
Independence, Comonotonicity, Countermonotonicity

Implicit Copulas:
Copulas extracted from common multivariate distributions; if the multivariate distributions are tractable, 
this might facilitate the simulation of the copulas
Special case: Gaussian copula, t copula, elliptical copula

Explicit copulas:
Copulas defined in terms of simple formulas
A prominent class are Archimedian Copulas.


Independence copula:
\(C(u_1, ..., u_d) = \prod_{i = 1}^d u_i \quad (u_1, ..., u_d) \in [0, 1]\)
Independence of the components of a random vector \(X = (X_1, ..., X_d)^\top\) corresponds to this copula.

Comonotonic copula:
\(C(u_1, ..., u_d) = \min \{u_1, ..., u_d\} \quad (u_1, ..., u_d) \in [0, 1]\)

* \(C\) is the joint distribution function of \((U, ..., U)^\top\) with \(U \sim U([0, 1])\)
* An intuitive characterization is the following: The random vector \(X = (X_1, ..., X_d)^\top\) is associated to the comonotonic copula, 
if \(X_i = T_i(X_1)\) for increasing functions \(T_i, i = 2, 3, ..., d\).

Countermonotonic copula (only for \(d = 2\)):
\(C(u_1, u_2) = \max \{u_1 + u_2 - 1, 0\} \quad (u_1, u_2) \in [0, 1]\)

The copula \(C\) is the distribution function of \((U, 1 - U)^\top\) with \(U \sim U([0, 1])\).
The bivariate random vector \(X = (X_1, X_2)^\top\) is associated to the countermonotonic copula, 
if \(X_2 = T(X_1)\) for some decreasing function \(T\).

These examples of fundamental copulas are associated with very particular 
and extreme forms of dependence.
